\section{Preliminaries}
In this section, we first introduce the colored Jones polynomial, and then we define a quantum dilogarithm, and variants of the logarithm and the dilogarithm.
We also describe some of their properties.
\par
For a knot $K$ in the three-sphere $S^3$, we denote by $J_N(K;q)$ the colored Jones polynomial of $K$ associated with the $N$-dimensional irreducible representation of the Lie algebra $\frak{sl}(2;\C)$ \cite{Kirby/Melvin:INVEM1991,Reshetikhin/Turaev:INVEM1991}.
We normalize it so that $J_N(U;q)=1$ for the unknot $U\subset S^3$.
Note that $J_2(K;q)$ is the original Jones polynomial \cite{Jones:BULAM31985}.
\par
Let $E$ be the figure-eight knot.
We use the following formula obtained by K.~Habiro \cite[P.~36 (1)]{Habiro:SURIK2000} and T.T.Q.~Le \cite[P.~129]{Le:TOPOA2003}.
See also \cite{Masbaum:ALGGT12003}.
\begin{equation}\label{eq:J_N_fig8}
\begin{split}
  J_N(E;q)
  &=
  \sum_{k=0}^{N-1}
  \prod_{l=1}^{k}\left(q^{(N-l)/2}-q^{-(N-l)/2}\right)\left(q^{(N+l)/2}-q^{(N+l)/2}\right)
  \\
  &=
  \sum_{k=0}^{N-1}
  q^{-kN}
  \prod_{l=1}^{k}\left(1-q^{N-l}\right)\left(1-q^{N+l}\right).
\end{split}
\end{equation}
\par
Next, we define functions $T_N(z)$, $\L_{0}(z)$, $\L_{1}(z)$, and $\L_{2}(z)$, which are related as
\begin{align*}
  T_{N}(z)
  &=
  \frac{N}{\xi}\L_{2}(z)+O(1/N)\quad(N\to\infty),
  \\
  \frac{d\,\L_{2}}{d\,z}(z)
  &=
  -2\pi\i\L_{1}(z),
  \\
  \frac{d\,\L_{1}}{d\,z}(z)
  &=
  -\L_{0}(z).
\end{align*}
\par
Let $\xi$ be a positive real number.
We define $C_{0}:=(-\infty,-1]\cup\{e^{\tau\i}\mid0\le\tau\le\pi\}\cup[1,\infty)$ and $C_{\theta}:=e^{\i\theta}C_{0}$, where $\theta$ is a positive real number with $\tan\theta<\frac{\pi}{\xi}$ and we orient $C_{0}$ from left to right.
\par
Consider the following integrals:
\begin{equation*}
  \int_{C_{\theta}}\frac{e^{(2z-1)x}}{x\sinh(x)\sinh(\gamma x)}\,dx,
  \quad
  \int_{C_{\theta}}\frac{e^{(2z-1)x}}{x^{m}\sinh(x)}\,dx,
\end{equation*}
where $m=0,1,2$, $\gamma:=\frac{\xi}{2N\pi\i}$ for an integer $N\ge2$, and $z$ is a complex number with $0<\Re(ze^{\i\theta})<\cos\theta$.
Here we follow \cite{Faddeev:LETMP1995} to introduce $T_N(z)$, which plays an important role in the paper.
Note that the set of the poles of the integrands are
\begin{equation*}
  \{k\pi\i\mid k\in\Z\}\cup\{2lN\pi^2/\xi\mid l\in\Z\}
  \quad\text{and}\quad
  \{k\pi\i\mid k\in\Z\},
\end{equation*}
respectively.
Therefore, if $N$ is large enough, then the only pole inside the unit circle centered at the origin is $0\in\C$ and so the path of integral $C_{\theta}$ avoids the poles.
Proofs of their convergences (Lemmas~\ref{lem:converge_T} and \ref{lem:integrals} below) are given in Section~\ref{sec:proofs_lemmas}.
\begin{lem}\label{lem:converge_T}
The integral $\int_{C_{\theta}}\frac{e^{(2z-1)x}}{x\sinh(x)\sinh(\gamma x)}\,dx$ converges if $z\in\C$ satisfies $-\frac{\xi\sin\theta}{4N\pi}<\Re(ze^{\theta\i})<\cos\theta+\frac{\xi\sin\theta}{4N\pi}$.
\end{lem}
\begin{lem}\label{lem:integrals}
If $z\in\C$ satisfies $0<\Re(ze^{\theta\i})<\cos\theta$, then the integral $\int_{C_{\theta}}\frac{e^{(2z-1)x}}{x^{m}\sinh(x)}\,dx$ converges for $m=0,1,2$.
\end{lem}
\par
Now define functions $T_N(z)$ and $\L_{m}(z)$ ($m=0,1,2$) by using the integrals above.
\begin{defn}
Fix an integer $N\ge2$ and put $\gamma:=\frac{\xi}{2N\pi\i}$.
We define
\begin{equation*}
  T_N(z)
  :=
  \frac{1}{4}
  \int_{C_{\theta}}
  \frac{e^{(2z-1)x}}{x\sinh(x)\sinh(\gamma x)}\,dx,
\end{equation*}
for a complex number $z$ with $-\frac{\xi\sin\theta}{4N\pi}<\Re(ze^{\theta\i})<\cos\theta+\frac{\xi\sin\theta}{4N\pi}$.
\end{defn}
\par
\begin{defn}
For a complex number $z$ with $0<\Re(ze^{\i\theta})<\cos\theta$, we define
\begin{align*}
  \mathcal{L}_0(z)
  &:=
  \int_{C_{\theta}}\frac{e^{(2z-1)x}}{\sinh(x)}\,dx,
  \\
  \L_1(z)
  &:=
  -\frac{1}{2}
  \int_{C_{\theta}}\frac{e^{(2z-1)x}}{x\sinh(x)}\,dx,
  \\[5mm]
  \L_2(z)
  &:=
  \frac{\pi\i}{2}
  \int_{C_{\theta}}\frac{e^{(2z-1)x}}{x^2\sinh(x)}\,dx.
\end{align*}
\end{defn}
\begin{lem}\label{lem:L0_1_2}
For $m=0,1,2$, we calculate $\L_{m}(z)$ as follows.
\begin{align}
  \L_0(z)
  &=
  \frac{-2\pi\i}{1-e^{-2\pi\i z}},
  \label{eq:L0}
  \\
  \L_1(z)
  &=
  \begin{cases}
    \log\left(1-e^{2\pi\i z}\right)&\text{if $\Im{z}\ge0$,}
    \\[3mm]
    \pi\i(2z-1)+\log\left(1-e^{-2\pi\i z}\right)&\text{if $\Im{z}<0$,}
  \end{cases}
  \label{eq:L1}
  \\[5mm]
  \L_2(z)
  &=
  \begin{cases}
    \Li_2\left(e^{2\pi\i z}\right)&\text{if $\Im{z}\ge0$,}
    \\[3mm]
    \pi^2\left(2z^2-2z+\frac{1}{3}\right)-\Li_2\left(e^{-2\pi\i z}\right)
    &\text{if $\Im{z}<0$}.
  \end{cases}
  \label{eq:L2}
\end{align}
\end{lem}
Proofs are also given in Section~\ref{sec:proofs_lemmas}.
\begin{rem}
In \eqref{eq:L1} and \eqref{eq:L2}, we use $\log(1-x)$ and $\Li_2(x)$ only for $|x|\le1$ ($x\ne1$).
\end{rem}
%%%%%%%%%%%%%%%%%%%%%%%%%%%%%%%%%%%%%%%%%%%%%%%%%%%%%%%%%%%%%%%%%%%%%%%%%%%%%%%%%%%%%%%%%
Since
\begin{equation*}
\begin{split}
  T_{N}(z-\gamma/2)-T_{N}(z+\gamma/2)
  &=
  \int_{C_{\theta}}
  \frac{e^{(2z-\gamma-1)x}-e^{(2z+\gamma-1)x}}{4x\sinh(x)\sinh(\gamma x)}\,dx
  \\
  &=
  -
  \int_{C_{\theta}}
  \frac{e^{(2z-1)x}}{2x\sinh(x)}\,dx
  \\
  &=
  \L_1(z),
\end{split}
\end{equation*}
we have the following corollary.
\begin{cor}
If $0<\Re(ze^{\theta\i})<\cos\theta$, then we have
\begin{equation}\label{eq:lem}
  \frac{\exp\left(T_{N}(z-\gamma/2)\right)}{\exp\left(T_{N}(z+\gamma/2)\right)}
  =
  1-e^{2\pi\i z}.
\end{equation}
\end{cor}
\par
We will show relations among the functions $T_{N}(z)$ and $\L_{m}(z)$ ($m=0,1,2$).
\par
We can prove that $\frac{1}{N}T_{N}(z)$ uniformly converges to $\frac{1}{\xi}\L_{2}(z)$.
More precisely, we have the following proposition.
See \cite[Proposition~A.1]{Ohtsuki:QT2016}.
\begin{prop}\label{prop:T_N_L_2}
For any positive real number $M$ and a sufficiently small positive real number $\nu$, we have
\begin{equation*}
  T_N(z)
  =
  \frac{N}{\xi}\L_2(z)
  +
  O(1/N)
  \quad(N\to\infty)
\end{equation*}
in the region
\begin{equation*}
  \{z\in\C\mid\nu\le\Re(ze^{\theta\i})\le\cos\theta-\nu,|\Im{z}|\le M\}.
\end{equation*}
In particular, the function $\frac{1}{N}T_N(z)$ uniformly converges to $\frac{1}{\xi}\L_2(z)$ in the region above.
\end{prop}
A proof is given in Section~\ref{sec:proofs_lemmas}.
Since we may regard $\L_2(z)$ as a variant of the dilogarithm from Lemma~\ref{lem:L0_1_2}, the function $T_{N}(z)$ is another quantum dilogarithm.
\par
The functions $\L_{0}(z)$, $\L_{1}(z)$, and $\L_{2}(z)$ are related as follows.
\begin{lem}\label{lem:der_L}
The derivatives of $\L_1(z)$ and $\L_2(z)$ are given as follows:
\begin{align*}
  \frac{d\,\L_2}{d\,z}(z)
  &=
  -2\pi\i\L_1(z),
  \\
  \dfrac{d\,\L_1}{d\,z}(z)
  &=
  -\L_0(z).
\end{align*}
\end{lem}
\begin{proof}
We have
\begin{equation*}
\begin{split}
  \frac{d\,\L_2}{d\,z}(z)
  &=
  \frac{\pi\i}{2}\int_{C_{\theta}}\left(\frac{d}{d\,z}\frac{e^{(2z-1)x}}{x^2\sinh(x)}\right)\,dx
  \\
  &=
  \pi\i\int_{C_{\theta}}\frac{e^{(2z-1)x}}{x\sinh(x)}\,dx
  \\
  &=
  -2\pi\i\L_1(z),
\end{split}
\end{equation*}
and
\begin{equation*}
\begin{split}
  \frac{d\,\L_1}{d\,z}(z)
  &=
  -\frac{1}{2}\int_{C_{\theta}}\left(\frac{d}{d\,z}\frac{e^{(2z-1)x}}{x\sinh(x)}\right)\,dx
  \\
  &=
  -\int_{C_{\theta}}\frac{e^{(2z-1)x}}{\sinh(x)}\,dx
  \\
  &=
  -\L_0(z),
\end{split}
\end{equation*}
completing the proof.
\end{proof}